\subsection{Experimental Plan and Architecture}
\label{sec:app1_experimental_plan}

% Rubric Requirements (15 Pts):
% 1. Experimental setting, training procedures, and robustness under limited computational resources.
% 2. Models used, architectural differences, and computed parameter counts.
% 3. Hyperparameter search techniques (grid search or better techniques).
% 4. Ablation table with different hyperparameters.
% 5. Metric selection rationale and additional metrics explored.
% 6. System effort required and scenario for the final viva demonstration.

\paragraph{Experimental Setting and Robustness}
% Describe training hardware constraints, precision settings, batching, optimizer choice, and techniques used to ensure stable results (e.g., cross-validation, multiple seeds).
[Describe the experimental environment, training configuration, and measures taken to ensure empirical robustness given compute bounds.]

\paragraph{Model Architectures and Parameter Counts}
% List the model variants used, provide exact parameter counts, and articulate the key architectural modifications that distinguish them.
[Detail model architectures, parameter counts, and structural modifications distinguishing this approach from baselines.]

\paragraph{Hyperparameter Optimization and Ablation Study}
% Describe the hyperparameter search method (e.g., grid search, Bayesian optimization) and summarize performance across configurations in Table~\ref{tab:app1_ablation}.
[Explain how optimal hyperparameters were determined across algorithms and baselines.]

\begin{table*}[t]
\centering
\small
\begin{tabular}{lcccc}
\toprule
Configuration & Learning Rate & Batch Size & Primary Metric & Secondary Metric \\
\midrule
Baseline & $1 \times 10^{-4}$ & 16 & 0.00 & 0.00 \\
Variant A & $5 \times 10^{-5}$ & 32 & 0.00 & 0.00 \\
Variant B & $2 \times 10^{-5}$ & 32 & 0.00 & 0.00 \\
Full Model & $2 \times 10^{-5}$ & 64 & 0.00 & 0.00 \\
\bottomrule
\end{tabular}
\caption{Hyperparameter ablation study for Approach 1.}
\label{tab:app1_ablation}
\end{table*}

\paragraph{Evaluation Metrics}
% Justify primary metrics chosen and discuss any exploratory or complementary metrics evaluated.
[Justify the primary evaluation metrics selected and describe any secondary or non-standard metrics tested.]

\paragraph{Viva Demonstration Scenario}
% Paint a clear, concrete walkthrough of what will be demonstrated interactively to course examiners during the viva.
[Provide an end-to-end scenario demonstrating user input, model inference, and output inspection during the viva session.]
